\notechapter{N-20260305}{Orthogonal Projection, Least Squares, and the Geometry of Approximation}

\section{The nearest point is found by a right angle}

Suppose a vector \(y\) does not lie in a subspace \(W\).  There are infinitely many vectors in \(W\), so what could make one of them the best approximation to \(y\)?

The answer is visible before it is algebraic.  If \(p\in W\) is the closest point, then the error
\[
  r=y-p
\]
must be perpendicular to every direction in \(W\).  Otherwise one could move a little inside \(W\) in a direction having positive inner product with \(r\), shortening the distance.

This gives the projection principle:
\[
  p=\operatorname{proj}_W y
  \qquad\Longleftrightarrow\qquad
  p\in W,\quad y-p\perp W.
\]
The right angle is not a decorative consequence.  It is the condition that singles out the nearest point.

\begin{figure}[H]
\centering
\includegraphics[width=0.92\textwidth,height=0.72\textheight,keepaspectratio]{figures/notes-md/N-20260305/N-20260305-fig-01-projection-onto-two-dimensional-subspace.png}
\caption{The best approximation is characterized by an orthogonal residual.  Moving within the subspace cannot reduce the remaining error.}
\end{figure}

\section{Projection onto one direction}

Let \(u\ne0\).  A candidate point on the line \(\operatorname{span}\{u\}\) has the form \(p=\alpha u\).  Orthogonality of the error says
\[
  \langle y-\alpha u,u\rangle=0.
\]
Hence
\[
  \alpha=\frac{\langle y,u\rangle}{\langle u,u\rangle},
  \qquad
  \operatorname{proj}_{u}y
  =\frac{\langle y,u\rangle}{\langle u,u\rangle}u.
\]
If \(u\) is a unit vector, this simplifies to
\[
  \operatorname{proj}_{u}y=\langle y,u\rangle u.
\]

Take
\[
  y=\begin{bmatrix}3\\1\end{bmatrix},
  \qquad
  u=\begin{bmatrix}1\\2\end{bmatrix}.
\]
Then
\[
  \langle y,u\rangle=5,
  \qquad
  \langle u,u\rangle=5,
\]
so
\[
  p=u=\begin{bmatrix}1\\2\end{bmatrix},
  \qquad
  r=y-p=\begin{bmatrix}2\\-1\end{bmatrix}.
\]
Indeed,
\[
  r^Tu=2-2=0.
\]
The decomposition
\[
  y=p+r
\]
separates the part visible inside the line from the part the line cannot represent.

\section{A basis turns the geometric condition into equations}

Let the columns of
\[
  B=\begin{bmatrix}b_1&\cdots&b_m\end{bmatrix}
  \in\mathbb R^{n\times m}
\]
form a basis of \(W\).  Every candidate approximation is \(Bc\) for some coefficient vector \(c\).  The residual condition
\[
  y-Bc\perp W
\]
means that it is orthogonal to each basis vector.  In matrix form,
\[
  B^T(y-Bc)=0.
\]
Therefore
\[
  B^TBc=B^Ty.
\]
If the columns of \(B\) are independent, the Gram matrix \(B^TB\) is positive definite and invertible, giving
\[
  c=(B^TB)^{-1}B^Ty
\]
and
\[
  \boxed{
  \operatorname{proj}_W y
  =B(B^TB)^{-1}B^Ty.}
\]

This formula looks more complicated than the one-dimensional version only because the basis vectors need not be orthogonal.  The Gram matrix records their mutual angles and corrects for the fact that the coordinate directions interfere with one another.

\section{The projection matrix remembers the subspace, not the basis}

Define
\[
  P=B(B^TB)^{-1}B^T.
\]
Then
\[
  Py=\operatorname{proj}_W y.
\]
Two identities reveal its geometry:
\[
  P^T=P,
  \qquad
  P^2=P.
\]
Symmetry says the projection is orthogonal.  Idempotence says that once a vector has been projected, projecting again does nothing.

Although the formula uses a basis, the matrix depends only on \(W\).  If \(C=BR\) is another basis matrix with \(R\) invertible, then
\[
  C(C^TC)^{-1}C^T
  =BR(R^TB^TBR)^{-1}R^TB^T
  =B(B^TB)^{-1}B^T.
\]
Coordinates changed, but the geometric operation did not.

There is also a complementary projection:
\[
  I-P.
\]
It sends \(y\) to the residual \(r=y-Py\), which lies in \(W^\perp\).  Thus
\[
  \mathbb R^n=W\oplus W^\perp,
  \qquad
  y=Py+(I-P)y.
\]

\section{Least squares is projection written as a fitting problem}

Now consider a system
\[
  Ax\approx b
\]
with more equations than unknowns.  Exact equality may be impossible because \(b\notin R(A)\), the column space of \(A\).  The least-squares problem asks for
\[
  \min_x\|Ax-b\|_2.
\]
But every vector \(Ax\) lies in \(R(A)\).  Therefore we are simply looking for the point in the column space closest to \(b\).

At the minimizer \(\widehat x\), the residual
\[
  r=b-A\widehat x
\]
must be orthogonal to every column of \(A\).  Hence
\[
  A^Tr=0,
\]
which gives the normal equations
\[
  \boxed{A^TA\widehat x=A^Tb.}
\]

The phrase ``normal equations'' comes from this normality: the residual is normal, or perpendicular, to the allowable model space.

\section{A complete line-fitting example}

Suppose three measured points are
\[
  (0,1),\qquad (1,2),\qquad (2,2).
\]
We fit a line
\[
  y\approx\alpha+\beta x.
\]
The design matrix and observation vector are
\[
  A=
  \begin{bmatrix}
    1&0\\
    1&1\\
    1&2
  \end{bmatrix},
  \qquad
  b=
  \begin{bmatrix}
    1\\2\\2
  \end{bmatrix}.
\]
The normal equations are
\[
  \begin{bmatrix}
    3&3\\
    3&5
  \end{bmatrix}
  \begin{bmatrix}
    \alpha\\\beta
  \end{bmatrix}
  =
  \begin{bmatrix}
    5\\6
  \end{bmatrix}.
\]
Solving gives
\[
  \widehat\alpha=\frac76,
  \qquad
  \widehat\beta=\frac12.
\]
Thus the fitted line is
\[
  y=\frac76+\frac12x.
\]
The predicted data are
\[
  A\widehat x=
  \begin{bmatrix}
    7/6\\5/3\\13/6
  \end{bmatrix},
\]
and the residual is
\[
  r=b-A\widehat x
  =
  \begin{bmatrix}
    -1/6\\1/3\\-1/6
  \end{bmatrix}.
\]
A direct check gives
\[
  A^Tr=
  \begin{bmatrix}
    -1/6+1/3-1/6\\
    0+1/3-2/6
  \end{bmatrix}
  =0.
\]
The fit is not exact, but no infinitesimal change in either the intercept or slope can improve the squared error to first order.

\section{The Pythagorean proof of optimality}

Let \(p=Py\) be the orthogonal projection of \(y\) onto \(W\).  For any other \(w\in W\), write
\[
  y-w=(y-p)+(p-w).
\]
The first term lies in \(W^\perp\), while the second lies in \(W\).  Therefore they are orthogonal, and Pythagoras gives
\[
  \|y-w\|^2
  =\|y-p\|^2+\|p-w\|^2.
\]
Hence
\[
  \|y-w\|\ge\|y-p\|,
\]
with equality only when \(w=p\).

This proof explains both existence and uniqueness in finite dimensions.  It also shows why the residual condition is stronger than a derivative calculation: the whole excess error is exactly the squared distance from \(w\) to the projection.

For least squares, the same identity becomes
\[
  \|b-Ax\|^2
  =\|b-A\widehat x\|^2
  +\|A(x-\widehat x)\|^2.
\]
Every competing parameter vector pays the irreducible residual plus an additional nonnegative model-space error.

\section{Orthonormal coordinates remove the Gram matrix}

If the columns of \(Q\) are orthonormal, then
\[
  Q^TQ=I
\]
and projection simplifies to
\[
  P=QQ^T.
\]
The coefficients are simply inner products:
\[
  c=Q^Ty.
\]
This is one reason orthonormal bases are so useful: geometry and coordinates agree without a correction matrix.

For a full-column-rank matrix \(A\), a QR factorization writes
\[
  A=QR,
\]
where \(Q^TQ=I\) and \(R\) is upper triangular.  Then
\[
  \|Ax-b\|_2
  =\|QRx-b\|_2.
\]
Multiplying by an orthogonal completion of \(Q\) separates the reachable and unreachable parts of \(b\), and the least-squares solution is obtained from
\[
  Rx=Q^Tb.
\]

This is computationally preferable to explicitly forming \(A^TA\).  The condition number is squared by the normal equations:
\[
  \kappa_2(A^TA)=\kappa_2(A)^2.
\]
Thus the normal equations are conceptually illuminating but can be numerically fragile.  QR preserves the projection geometry without amplifying the conditioning in this way.

\section{Rank deficiency changes uniqueness, not the geometry}

If the columns of \(A\) are dependent, the fitted vector \(A\widehat x\) can still be unique even though the coefficient vector is not.  Indeed, all least-squares minimizers produce the same projection
\[
  A\widehat x=Pb,
\]
but two parameter vectors may differ by an element of the null space:
\[
  x_1-x_2\in N(A).
\]

This distinction is easy to miss in applications.  The model prediction can be completely determined while the internal coefficients remain ambiguous.  Collinear features in regression are a familiar example.

The Moore--Penrose inverse, studied later in more detail, chooses the minimizer orthogonal to \(N(A)\).  It is therefore the least-squares solution with minimum Euclidean norm:
\[
  \widehat x=A^+b.
\]
Here the pseudoinverse is a canonical coordinate choice, while the geometric heart of the problem remains projection onto \(R(A)\).

\section{Weighted least squares changes the notion of perpendicular}

Sometimes residual coordinates should not be treated equally.  If \(W\) is symmetric positive definite, define
\[
  \langle u,v\rangle_W=u^TWv,
  \qquad
  \|u\|_W^2=u^TWu.
\]
Weighted least squares minimizes
\[
  \|Ax-b\|_W^2=(Ax-b)^TW(Ax-b).
\]
Differentiation or weighted orthogonality gives
\[
  A^TW(A\widehat x-b)=0,
\]
so
\[
  A^TWA\widehat x=A^TWb.
\]

The algebra resembles ordinary least squares, but the geometry has changed.  The residual is no longer Euclidean-orthogonal to the column space; it is orthogonal in the \(W\)-inner product.  Equivalently, after writing \(W=C^TC\), one solves the ordinary least-squares problem
\[
  \min_x\|CAx-Cb\|_2.
\]
This is how measurement confidence, covariance information, and anisotropic geometry enter the same projection framework.

\section{Projection beyond finite-dimensional matrices}

The finite-dimensional picture survives in Hilbert spaces.  If \(H\) is a Hilbert space and \(M\subseteq H\) is a closed subspace, then every \(y\in H\) has a unique decomposition
\[
  y=p+r,
  \qquad
  p\in M,\quad r\in M^\perp.
\]
The point \(p\) is the unique best approximation to \(y\) from \(M\).

The word ``closed'' matters.  Without it, an infimum distance may exist without being attained inside the subspace.  Completeness and closedness are the infinite-dimensional replacements for the compact geometric intuition that works automatically in finite dimensions.

Fourier series are a central example.  Truncating a Fourier expansion is orthogonal projection onto a finite-dimensional space of trigonometric modes.  Polynomial approximation, finite elements, and Galerkin methods use the same principle with different trial spaces.  The matrix normal equations are therefore one coordinate shadow of a much broader approximation theorem.

\section{What should remain after the formulas}

Projection and least squares are one story.  A model can reproduce only vectors in its column space.  The best approximation is obtained by dropping a perpendicular to that space.  A basis converts the perpendicularity condition into normal equations; an orthonormal basis makes the coordinates transparent; QR turns that geometry into a stable algorithm.  Rank deficiency separates uniqueness of predictions from uniqueness of parameters, and weighted inner products alter what ``nearest'' means without changing the logical structure.

The enduring test is simple:
\[
  \text{best approximation}
  \quad\Longleftrightarrow\quad
  \text{residual orthogonal to every allowed correction}.
\]
That sentence is the bridge from a diagram with a right angle to regression, Fourier analysis, numerical PDEs, and Hilbert-space approximation.
